\chapter{Executive Overview}
\label{ch:executive_overview}


\section{Mission Objective and Value Proposition}

The High-speed Unmanned General-purpose Observatory (HUGO) is a reusable, modular and adaptable 50 kg transonic flying laboratory designed to bridge the gap between wind-tunnel testing and full-scale flight testing. Developed to accelerate the validation of next-generation aerospace technologies, HUGO provides a modular, low-cost platform for testing high-aspect-ratio (HAR) wings \cite{AFONSO201740, MA2024100983}, Active Flutter Suppression (AFS) \cite{Gilyard1983, Takarics2019}, and Gust Load Alleviation (GLA) systems \cite{Sanghi2022} in real-world atmospheric conditions.

Currently, researchers and industry professionals face significant financial and technical barriers when testing at Mach 0.75. Traditional transonic wind tunnels suffer from wall-interference and scaling limitations \cite{Stenfelt2015}, often costing upwards of €40,000 per occupancy day \footnote{\label{foot:schrijer-exec}In personal communication with Ferdinand Schrijer on 04.05.2026.}. Conversely, manned flight testing introduces severe safety risks and requires lengthy, expensive certification processes \cite{Spitz2001}. 

Project HUGO disrupts this paradigm by offering a certified, reconfigurable testbed that operates on Sustainable Aviation Fuel (SAF) to minimize altitude emissions \cite{Gaillot2023}. At a targeted commercial rate of €25,000 per flight day, HUGO significantly reduces the programmatic costs of aeroelastic research while delivering high-fidelity, free-flight data.

To evaluate the commercial feasibility of HUGO, a market volume analysis was conducted. While the Total Addressable Market (TAM) encompasses the broader global expenditure on aerospace testing and validation, HUGO's initial beachhead strategy specifically targets the domestic Dutch testing ecosystem. Based on the collective testing budgets of major institutions such as the Netherlands Aerospace Centre (NLR), the German-Dutch Wind Tunnels (DNW), and TU Delft, the Serviceable Available Market (SAM) is valued at €71 million \todo{cite mid term report} annually. By providing a direct alternative to the high-speed aeroelastic testing traditionally performed in wind tunnels, HUGO aims to capture roughly 20\% \footref{foot:schrijer-exec} of this regional market, yielding a Serviceable Obtainable Market (SOM) of €14.2 million per year. Furthermore, the drastic reduction in testing costs and lead times is expected to induce latent demand, expanding the overall market boundaries by allowing smaller academic and commercial entities to conduct research that was previously cost-prohibitive.

\section{Final Configuration and Modularity}

Following a rigorous Analytic Hierarchy Process (AHP) trade-off, HUG-CFG-303 was selected as the optimal configuration. This high-wing configuration features static wing mounting points, a T-tail, and a modular canard port, providing maximum experimental versatility without the structural penalties associated with more complex airframes with multiple wing mounting points. 

To fulfill its role as a universal testbed, HUGO is designed around a defined ``Planform Envelope.'' This allows clients to test wings with aspect ratios ranging from 5 to 27, sweeps up to 20$^{\circ}$, and various dihedral and thickness profiles, all while ensuring stable and controllable flight characteristics \cite{nelson1998flight}. For missions that do not require a custom client wing, HUGO employs a Standard Planform featuring a NASA SC(2)-0012 supercritical airfoil \cite{Harris1990}, optimized for stable, high-subsonic operations. This is paired with a standard HUGO empennage, which is a T-tail. 

Additionally, HUGO features a canard mounting point, which is covered to prevent an aerodynamic penalty when not being used. This allows HUGO to be configured with a canard on client request, in order to be able to test aircraft configurations which feature a canard. An additional benefit of this feature is that it allows for the testing of highly experimental planform at a smaller scale on these ports, without having to design, manufacture and test a new main planform.

\section{Engineering Design and Performance Outcomes}

The technical execution of HUGO ensures that the aircraft can perform demanding test profiles while adhering to the strict user and operational requirements:

\begin{itemize}
    \item \textbf{Structural Integrity:} The primary structure utilizes out-of-autoclave CYCOM 5320-1 carbon fiber epoxy prepreg \cite{CENTEA2015132}. The fuselage features a Nomex honeycomb core \cite{Boukharouba2014}, while the lifting surfaces utilize a Polymethacrylimide (PMI) foam core. Structural analysis confirms the airframe withstands the ultimate load factor of 9g, with a maximum normal bending stress of 696 MPa—well below the material's yield limit. The airframe has also been analyzed to withstand the handling and abuse loads involved with the operations of HUGO. Furthermore, aeroelastic modelling confirms flutter and static divergence occur well above the Mach 0.75 operating speed \cite{wright2008introduction}.
    \item \textbf{Propulsion and Power:} Thrust is provided by twin fuselage-mounted JetCat P100-RX-BL micro-turbines, fed by a dynamic fuel management system utilizing two 8.4L fuel bladders. This allows us to dynamically shift the CG during flight in orfer to change the stability of the aircraft to run different experiments. Electrical power is supplied by a custom 8S5P Molicel INR21700 battery pack \cite{Tsung-Hsun2026, Ganjeizadeh2014}, delivering 936 W to sustain the entire avionics, sensor, and propulsion network for the maximum 45-minute flight window.
    \item \textbf{Performance:} Payload-range analysis validates that HUGO can successfully complete its 30-minute nominal mission—including a 5-minute Mach 0.75 test segment at FL270, while holding sufficient fuel reserves for loitering and up to four missed approach attempts \cite{Koutz1952Loitering}.
\end{itemize}

\section{Operations, Certification, and Safety (RAMS)}

Operations and Logistics are a crucial part in the development of a successful mission. Together they cover all aspects of HUGO's day-to-day usage. For a mission to be successful, HUGO must operate in a controlled airspace, fly over 100km in a straight line, and to perform a test at Mach 0.75 for at leasts 5 minutes. Given its classification as a High-Risk Unmanned Aircraft System under EASA Special Conditions \cite{easa_light_uas_crd_is3}, HUGO's operational and safety framework is designed for maximum reliability. Operations are based at Den Helder Airport (EHKD), capitalizing on its 1.27 km runway, emergency services, and direct access to ATC-controlled airspace over the North Sea \cite{DenHelderAirport2013DHANieuwsbrief}. 

To facilitate the operations and reduce transport costs, HUGO is to be stored in a hangar in the area adjacent to Den Helder airport. Although not a commercial airport, Den Helder is the main gateway to off-shore platforms in the North Sea, thus daily traffic is expected. As such, the test flight plans will be coordinated with De Kooy Air Traffic Control. HUGO's propulsion system uses Sustainable Aviation Fuel (SAF), which will need to be stored in the hangar and brought to the airport for each flight day as Den Helder airport does not offer SAF refuelling services. Lastly, the handling of HUGO; transport, flight, and ground support will be performed by trained personnel.

In consultation with, and with guidance from EASA, HUGO was classified under SC-Light-UAS High risk. Successful certification of an Unmanned Aerial Vehicle (UAV) with such a high degree of modularity would require constant and expensive recertification, making this project infeasible. However, through the simultaneous application for a type certificate and a Design Organisation Approval (DOA), HUGO will be allowed to undergo major changes without requiring full recertification. Though more versatile, it is a process that will require an initial investment of \euro{10,170} for the certification, and then a yearly fee of \euro{11,050} for the maintenance of the Type Certificate. For compliance, a team of engineers will therefore have to apply for the DOA, and be responsible for approving any major changes proposed by a client.

Safety is ensured through a selective redundancy philosophy and a rigorous condition-based maintenance (CBM) or time-based maintenance (TBM) schedule. Critical systems, including the VECTOR-600 autopilot and power distribution, are monitored continuously \cite{Imani2022}. In the event of communication loss or envelope exceedance, autonomous stabilization and flight termination protocols guarantee safe recovery. The overall system reliability is calculated to ensure the initial fleet of five aircraft can easily sustain 100 operational test days per year \cite{Mencik2016, Lanzafame2024}.

\section{Sustainability Integration}








\begin{comment}

-Please Enter your section overview here: Summary of outcomes and their impact. Not too much explanation of procedure and how you got to the outcomes.

:Outline of Design Activities:

- analysis is split into mostly independent fixed group, planform envelope and standard planform analyses. 

- Fixed is the sizing of lg, fuselage and its interior and engines, then it is checked against housing all compatible planforms in the planform envelope analysis. 

- Further examination is conducted on the standard planform, the planform that is provided to the customer if they do not bring their own wing. 

- The analyses are verified by a combination of comparisons with reference, cross-verification with flow 5 and sensitivity studies, as well as hand calculations and sanity checks.

-Planning:


-Requirements:
    - about planform envelope, mention that we have defined the parameter ranges and the most constraining combinations to evaluate against


-Market Analysis:


-Technical risk assessment and certification:
    - operations and logistics: similar content to the midterm. A final risk matrix was constructed accounting for all the development and operational risks. Preventive measures and contingency scenarios were developed for each risk.
    - certification: similar content to the midterm, no new findings were made.

-Operations and logistics:
    - similar content to the midterm, except that now we also define exemplary missions that can be performed by HUGO.


-V and V: Talked with Hania initially, ad we may merge it with Chapter 3, ignore 4 now - Marek


-Reliability, Availability, Maintainability and Safety:


-Configuration layout and system characteristics:

We identified the resources first and budgets, outlined the factors that influence them and the most important parameters that have to be taken into account. We could not provide numbers for the budget at the begining outside the requirements that did not provide numnerical values for parts such as electronics stability and so on. 
Moving to internal configuration: 
    - electronic system has been configured, sized and proportioned in accordance with the mission necessities
    - we have an autopilot that can perform the majority of the aerodynamic measurments accompanied by a flight companion that performs calculations and collects data from all the other sensors
    -we store the data and tranmit it as well usint VHF waves
    -for the power budget we build a battery that uses 40 individual cells to supply 936 W for 45 minutes covering the flight and the additional 3 landing attempts. 
    -the power system takes into account the presence of the sensors, while the sensors use 51.7 W, the electronic components unse 841 W, for a combined 892.7, which is within the range of the battery.
    -propulsion system was also sized, having all the components brought toghther and analysed 
    -the fuel was computed and integrated into the numbers, the fuel subsystem utilising 2 fuel bags that amount to 10L of fuel for the mission
    -the fule bagas are used for the CG translation experiment where we move the CG from a stable configuration to an unstable one and that is why are are using larger fuel tanks, and the fuel tanks can be movable inside the fuselage depending on the platform that potential clients may bring  
    -we are using the 2 Jetcat PBJ 100 engines that are mounted to the aft section of the fuselage

Sensors-Carl

Landing gear-Gui

- Landing gear was sized for tip back angle, turn over angle, wing strike angle, nose landing gear controlability, landing gear strut length, and main landing gear horizontal track. 

- Optimised using scipy optimise for an ideal cylyndrical fuselage based on the dimensions from FlexOp

- Landing gear was then implemented in CAD with non idealised fuselage, and the constraints checked iteratively until all were met.

- The mechanism was designed to allow for the meeting of the previously mentioned constraints, and to stow away teh landing gear for flight. (might be nice to put the picture from the LG section here).

- Landing gear was simulated as a spring and dashpot in parallel to absorb landing shocks using the python control library. Main constraint was to minimise landing gear vertical displacement (compression) due to downward velocity at landing, and prevent over-damping as to dissipate the applied force over a longer period of time. A final choice of C = 781.115 and K = 48168.2, for a maximum allowable displacement of 6.84 cm at a downward velocity of 2m/s with an added safety factor of 1.5 as prescribed by EASA in SC-UAS High Risk. The effect of transverse loads on the landing gear during landing were taken into account during the material selection.
    

-Material characteristics + MAI (they are now one section):
    Materials were selected to ensure reliable performance in the high-subsonic regime, including resistance to aeroelastic effects such as flutter, while keeping manufacturing costs low for limited production volumes.

    -Steel alloy was chosen for the landing gear and engine pylon due to its high melting temperature and ability to achieve the required strength with smaller strut diameters. Aluminium alloy was used for the connection ports, where loads and thermal demands are comparatively low.

    -A sandwich composite panel was selected for the fuselage and tailcone, using a prepreg system suitable for out-of-autoclave (OOA) curing. This reduces manufacturing cost and cycle time, making it appropriate for low-rate production. Although aluminium could offer higher intrinsic strength, the aramid core avoids galvanic corrosion and provides sufficient compliance to conform to fuselage curvature without buckling. The 48 kg core variant was selected for adequate compressive strength to resist local crushing during cure.

    -For the wings and tail stabilisers, a PMI foam core was used due to its cost-effectiveness and suitable mechanical performance for low-temperature curing cycles.

    Once the material has been selected, the next step is to come up with the MAI plan. The manufacturing and assembly of the plan are just as important as the material selection for the final product quality and material characteristics.

    - The selected prepreg is optimised for OOA/VBO curing cycle. A continuous curing cycle was selected, during which the part remains on the mould. It also provided the lowest manufacturing time and the best dimensional accuracy. The honeycomb fuselage and tailcone are manufactured by first curing the outerskins in a separate curing cycle to achieve the desired aerodynamic finish, then the honeycomb core and inner skin are layered on each half with films of adhesive and the final assembly (both moulds held together undergoes a final curing cycle. The manufacturing of the foam wings consists of CNC machining the foam and preparing it for the layering of the prepreg with adhesive film. It then undergoes the same curing cycle. The landing gear and engine pylons are CNC-machined, heat-treated, and undergo an application of a protective surface coating.

    -The assembly process for the honeycomb structures consists of connecting the two mould halves that would then undergo a final curing cycle to produce the monocoque structure. After the outer skins have been cured and the honeycomb core and inner skin have been layered on them, the outer skins have to be connected by an adhesive film by overlapping their edges as a lap joint. An adhesive film layer is also applied at the honeycomb interfaces, essentially connecting the two mould halves, and the assembly undergoes the final curing cycle.

    -the pylon to fuselage connections will utilise metal inserts that will be inserted in holes drilled in the sandwich structure. The exposed honeycomb core will be sealed at those locations with an epoxy that bonds the metal inserts to the inside of the fuselage structure. Titanium is used due to its improved electrochemical compatibility with carbon fiber composites.

    -the landing gear assembly uses bolted connections for the rigid interfaces and pin joints where relative motion is required

    The integration is the same as for the midterm. Accessibility and system dependencies are driving. 

-Structural analysis:
Wing structure
    -we wrote code for the wing that tests its limitations in terms of structural rigidity and flutter analysis, as they were the driving parameters for the sizing iteration process
    -the wing was tested to meet the deflection, buckling requirements and to ensure that all the stresses do not go over the material limits such as shear, normal and torsional stresses
    - the deflection and buckling req are validated for a load factor as 6 as stated buy the regulations, while the structural characteristics are tested for a load factor of 9.


-Perfromance analysis:


-Stability and control:


%-Manufacturing, assembly and integration:


-Susteainability:


-Config and layout for different mission  concepts:

\end{comment}